\section{What remains, and the structure of the problem}\label{sec:remaining}

\subsection{The obligations}
\begin{obligation}[\cid{O.1}, FIN]
$\Sha(E/\Q)$ is finite. Equivalently, every primary component is finite and $\Sha[p]=0$ for all but finitely many $p$.
\end{obligation}
\begin{obligation}[\cid{O.2}, ORDER9]
If $\Sha$ is finite, then $\#\Sha=9$. Given FIN, it suffices to show $\#\Sha<2601$. The bound $\#\Sha[n]<2601$ for all $n$ is equivalent to FIN \emph{together with} ORDER9 (Proposition~\ref{prop:uniform}), not to this conditional assertion alone.
\end{obligation}
\begin{obligation}[\cid{O.3}, Q36]
$q=36$. No comparison between a $p$-adic leading term and the real regulator is available. (This is narrower than saying that $p$-adic methods cannot give $q$.) Even a proof of $q=36$ would leave ORDER9.
\end{obligation}

Together these three obligations are equivalent to full BSD for $E$ (Theorem~\ref{thm:endpoint}): the group obligation is finiteness together with order nine, and the real obligation is $q=36$. FIN does not imply ORDER9; given FIN, the settled components allow $\#\Sha=9t^2$ with every prime factor of $t$ at least $53$. Every other statement in this paper is settled at its stated grade (Appendix~\ref{app:grades}), so these endpoint obligations form the current known frontier. The claim index also keeps \cid{B.8} (the unevaluated witness at $4547$, a regime within the prime-by-prime strategy) and \cid{O.4} (Question~\ref{q:uniform}, a proposed mechanism) open; neither is a further endpoint obligation. The frontier describes the endpoint; it is not a count of missing lemmas. In this revision:
\begin{itemize}[leftmargin=*]
\item the rank, the basis and the primary components at $p\le13$ (\cid{R.1}, \cid{R.4}, \cid{S.2}, \cid{S.3}, \cid{S.5}--\cid{S.13} and \cid{S.P}) do not depend on an unrefereed preprint; the components at $17\le p\le47$ (\cid{S.17-47}) apply Ota's theorem as stated in its arXiv version, and the conditional route results \cid{B.9}, \cid{B.13}, \cid{B.17}, \cid{B.18} and \cid{B.19} keep their explicit preprint dependencies;
\item the deductions of the primary components at $3,5,7,11,13$ are Lean-checked from their listed hypotheses (Remarks~\ref{rem:a4lean} and~\ref{rem:a5lean});
\item the $2$-primary component is an exact computation;
\item finiteness of $\Sha[3^\infty]$ and the lower bound $\Sha[3]\supseteq(\Z/3)^2$ do not depend on the main conjecture (Theorems~\ref{thm:sha3fin} and~\ref{thm:sha3lower}); the exact order $9$ still does (Remark~\ref{rem:su}).
\end{itemize}
Under the retained prime-by-prime strategy, FIN and ORDER9 need every good prime $p\ge53$ and the bad primes $2243$, $4547$ and $5843$ (the table below). Prime-by-prime work cannot finish on its own, and a proof need not use that strategy. For Q36, rationality of $q$ with denominator at most $174$ would already give $q=36$ (Proposition~\ref{prop:ratl}); bounded algebraicity (Proposition~\ref{prop:alg}) and the common certificate of Proposition~\ref{prop:commonB} are other sufficient routes. None of these antecedents is proved.

The primes not yet settled fall into five regimes:

\begin{center}
\begin{tabular}{@{}p{0.3\textwidth}p{0.62\textwidth}@{}}
\toprule
regime & available sufficient criteria\\ \midrule
good ordinary $p\ge53$ & the singleton criterion of \S\ref{subsec:singletons} and Criterion~\ref{crit:cyclo}, computable prime by prime; Theorem~\ref{thm:ordbridge} and Corollary~\ref{cor:unified}, not evaluable\\
good supersingular $p\ge53$ ($131,167,\dots$; infinitely many \cite{elkies}) & the singleton criterion and Criterion~\ref{crit:ss}, computable prime by prime; Theorem~\ref{thm:a3}, not evaluable\\
$2243$ & Theorem~\ref{thm:2243}: one coefficient, uncomputed; Criterion~\ref{crit:bh}, unevaluated\\
$5843$ & Criterion~\ref{crit:5843}: one coefficient, uncomputed; Criterion~\ref{crit:bh}, unevaluated\\
$4547$ & the retained two-prime Kurihara sufficient witness and Criterion~\ref{crit:bh}, both unevaluated\\ \bottomrule
\end{tabular}
\end{center}
Every criterion in the table works at a single prime. FIN would follow from vanishing outside a finite set, together with finiteness at each remaining prime. Separate theorems may handle ordinary, supersingular and bad primes; one theorem need not cover every regime. Finite prime support alone does not control the depths at the exceptional primes.

Targeted screens of the literature, with the screened sources recorded with the evidence, found no published theorem that supplies a uniform-in-$p$ input of either kind for any elliptic curve of rank at least two: neither nondegeneracy of the $p$-adic height pairing for almost all $p$, nor nonvanishing modulo $p$ of depth-two Mazur--Tate coefficients uniformly in $p$. Such inputs are known one pair $(E,p)$ at a time, by rigorous computation, as in \S\ref{sec:primary}. This describes the current known frontier of the sources read; it is not a claim that no such theorem exists.

\subsection{The structure of the problem}
The following observations come from reading the record as a whole. Items marked \emph{(builder)} are ours and are not in the cited records.

\begin{enumerate}[label=(\alph*),leftmargin=*]
\item \emph{One shared import (builder).} Skinner's multiplicative theorem supports three results at once: the only nonzero primary component (Theorem~\ref{thm:sha3}) and both multiplicative bad-prime routes (Theorem~\ref{thm:2243}, Criterion~\ref{crit:5843}). An actual defect in the imported theorem would affect these chains together. The failure of the project's attempted internal reconstruction does not establish such a defect. The good-prime components rest on Kato's divisibility instead, which is independent of it.

\item \emph{One computational trust base (builder).} Every analytic certificate is an exponentially weighted twisted $L$-series evaluated from overlapping coefficient corpora sharing one verifier lineage. Independent replays share the raw coefficients. A review of the integer coefficient verifier is therefore the single most valuable audit on the computational side.

\item \emph{Two lanes with a similar unresolved requirement.} The Stark lane over $\Q(\sqrt{-23})$ and the definite-theta lane over $\Q(\sqrt{-7})$ both lack an arithmetic evaluation of their specified witness; no identification of those witnesses has been proved. Neither provides a prime-independent integral object whose reductions control the $p$-adic witness. In analytic rank one that object is the Heegner index of a conductor-one Heegner point. For this curve the corresponding index over $\Q(\sqrt{-23})$ is infinite, since the rank there is three.
\end{enumerate}

\begin{question}[\cid{O.4}]\label{q:uniform}
Is some sufficient witness for this curve the specialization of a single integral object that does not depend on $p$, and whose nonvanishing or defect is controlled by one fixed integer? Candidate witnesses include $M_2(p)$, $I_3(p)$, $C_p$ and $h_p$. A positive answer that also covers the supersingular and bad primes would give finite prime support. Lemma~\ref{lem:minors} shows that, once such an integral model exists, the effective constant follows formally.
\end{question}

\subsection{Is this a good place to start?}
The assessment in this subsection is ours and is not in the cited records.

\emph{In favour.}
\begin{itemize}[leftmargin=*]
\item The rank and analytic rank are certified.
\item There is an explicit basis with a certified regulator.
\item $q$ is certified within $0.006$ of $36$.
\item The components at $p\le47$ are settled, and the only nonzero one, $(\Z/3)^2$, is explicit.
\item Certified auxiliary quadratic fields exist in rank two and in rank three.
\item A conditional criterion, with a Lean-checked core, covers every good prime $p\ge5$ with $p\nmid7N$.
\item Many mechanisms have been tested, and their failure points are recorded with scope.
\end{itemize}

\emph{Against.}
\begin{itemize}[leftmargin=*]
\item The conductor $6.5\cdot10^{15}$ makes naive full modular-symbol or quaternionic spaces have scale $10^{14}$--$10^{16}$.
\item The conductor is not squarefree, which excludes several published theorems in their printed form.
\item The additive, potentially supersingular prime $4547$ has no executable certificate in the retained implementation.
\item The split prime $2243$ carries an exceptional zero.
\item The exact order of the nonzero primary component uses Skinner's printed theorem, one step of whose printed proof is not established for this curve (Remark~\ref{rem:su}); its finiteness and its lower bound do not.
\end{itemize}

We leave to experts the judgment whether this balance makes $E$ a useful test case for rank-two methods. This paper and its repository are organized to make that judgment possible.
